% Section: P5P8-SYNTH thirteen-domain density matrix — D13 per-prompt reward stability (iter 164 JOB B)
% Fresh vein. Extends iter-160's twelve-domain matrix (D1-D12) by adding
% D13 = per-(method × step × prompt) reward-stability density on the N2
% reward tensors (2560 cells = 4 × 40 × 16). H1-H4 4/4 PASS — sharpest
% finding: per-prompt reward stability is structurally unreachable at G=8.

\subsubsection{D13: per-prompt reward stability extends the 12-domain matrix to 13}
\label{sec:synth-iter164-d13}

Iter-160 added D12 = P5 N2 step-aggregate reward-stability density on
the 160-cell N2 reward tensor (4 methods $\times$ 40 steps;
\secref{sec:synth-iter160}). The natural follow-up is the per-prompt
granularity: D13 on $4 \times 40 \times 16 = 2560$ cells. The motivation
is the GRPO training loop: each step sees 16 prompts $\times$ 8 rollouts
of binary reward. If per-prompt reward is well-measured at G=8, group
baselines are stable; if not, the per-prompt CI is dominated by the
binomial $n=8$ floor and reward stability is structurally unreachable.

\paragraph{Setup.}
For each (method $\times$ step $\times$ prompt) cell, we compute the
Wilson 95\% CI half-width on the 8-rollout binary reward proportion
and define $D_{13}(\mathrm{cell}, \varepsilon) = \mathbb{1}[\mathrm{half\text{-}width} < \varepsilon]$.
Density = (\#stable cells) / 2560.
$\varepsilon \in \{0.025, 0.05, 0.10\}$.

\paragraph{Headline findings (4/4 PASS, all structural).}
\begin{enumerate}
\item \textbf{$D_{13}@0.05$ density = 0.0000.} H1 PASS. None of the
2{,}560 (method $\times$ step $\times$ prompt) cells clears the
$\varepsilon = 0.05$ bar; the layer assignment is LOW.
\item \textbf{Wilson CI half-width floor = 0.1622 at $n=8$.} H2 PASS.
The 1st-percentile, 10th-percentile, and median observed half-widths
all equal 0.1622 (the theoretical floor when $k=0$ or $k=8$). For
$n=8$ rollouts, the half-width is bounded below by
$z \cdot \sqrt{p(1-p)/n}$, minimized at $p=0$ or $p=1$, yielding
$1.96 / (2 \cdot \sqrt{8}) \cdot \sqrt{(0 + 1.96^2/(4 \cdot 8)) / 1.48}$
$\approx 0.162$. Every $\varepsilon \leq 0.10$ is structurally infeasible.
\item \textbf{$D_{13} \ll D_{12}$.} H3 PASS. $D_{13} = 0$ vs
$D_{12} = 0.175$; ratio = 0. The per-prompt-vs-step granularity gap
is exactly the binomial $n=8$ floor.
\item \textbf{All four methods in LOW layer.} H4 PASS. grpo=0, aero=0,
gift=0, areal=0 (each on 640 cells). The structural floor binds
identically across the GRPO family.
\end{enumerate}

\paragraph{Operational implication for GRPO-family training.}
\begin{itemize}
\item \textbf{Per-prompt reward stability at $G=8$ is NOT measurable.}
Reporting per-prompt confidence intervals on a $G=8$ GRPO rollout is
structurally meaningless.
\item \textbf{To achieve $\varepsilon \leq 0.10$ per-prompt stability,
you need $G \geq 67$.} At $p=0.5$, the half-width is
$z \cdot \sqrt{0.25/n} \leq 0.10 \Rightarrow n \geq 96$. For
$\varepsilon \leq 0.05$ you need $G \geq 384$.
\item \textbf{Step-aggregate ($D_{12}$) is the canonical reporting
granularity at $G=8$.} $D_{12}@0.05 = 0.175$ with Wilson CI
$[0.124, 0.241]$ sits in MID — step-aggregate reward stability is a
meaningful metric, but per-prompt is not.
\end{itemize}

\paragraph{13-domain layer assignments.}
\textbf{LOW} = $\{D_1, D_6, D_7, D_{13}\}$ (4 domains; $D_{13}$ is the
first structural-LOW domain — its LOW status is determined by the
binomial $n=8$ floor, not by empirical sample sparsity).
\textbf{MID} = $\{D_2, D_3, D_4, D_8, D_9, D_{11}, D_{12}\}$ (7 domains).
\textbf{HIGH} = $\{D_5, D_{10}\}$ (2 domains).

\paragraph{Cross-paper coupling.}
\begin{itemize}
\item \textbf{P5P8-SYNTH iter-160 row 175 ($D_{12}$ step-aggregate):}
$D_{12}@0.05 = 0.175 \in$ MID; iter-164 $D_{13}$ at per-prompt = 0
forces LOW — the granularity gap is structural.
\item \textbf{P5 iter-141 (N10 reward tensor):} iter-141 produced the
160-cell N2 tensor; iter-164 reads the same tensor at finer granularity.
\item \textbf{P7 iter-163 row 177 (step-aggregate Pareto):} P7 controller
evaluation is at step granularity, consistent with iter-164's
recommendation that step-aggregate is the canonical reporting granularity
at $G=8$.
\end{itemize}